% Generated draft: requires mathematical and source review.
\documentclass[11pt]{article}
\usepackage[margin=25mm]{geometry}
\usepackage{fontspec}
\setmainfont{Latin Modern Roman}
\usepackage{amsmath,amssymb,mathtools,unicode-math}
\setmathfont{Latin Modern Math}
\usepackage{enumitem,needspace,xurl,hyperref}
\setlength{\parindent}{0pt}
\setlength{\parskip}{5pt}
\setlength{\emergencystretch}{3em}
\newcommand{\N}{\mathbb{N}}
\newcommand{\Z}{\mathbb{Z}}
\newcommand{\Q}{\mathbb{Q}}
\newcommand{\R}{\mathbb{R}}
\newcommand{\C}{\mathbb{C}}
\newcommand{\F}{\mathbb{F}}
\newcommand{\A}{\mathbb{A}}
\newcommand{\PP}{\mathbb{P}}
\newcommand{\E}{\mathbb{E}}
\newcommand{\eps}{\varepsilon}
\newcommand{\dd}{\,\mathrm{d}}
\begin{document}
\section*{UnsolvedMath ID: 10400210}
\textbf{Source problem code:} catalogue entry
\subsection{Definitions and mathematical statement}
In the usual unframed link setting, a link is a finite disjoint union of embedded circles in \(S^3\), up to ambient isotopy. A trivial link is an unlink: its components bound pairwise disjoint embedded disks. A standard \(3\)-move replaces inside a ball a trivial two-strand tangle by three consecutive half-twists, or performs the inverse replacement, leaving the link outside the ball fixed. The precise local picture and orientation convention must agree with the cited source.

Conjecture 12.4, the \(3\)-move conjecture attributed to Y. Nakanishi [305], asserts:
\[
 \forall\text{ links }L\ \exists\text{ an unlink }U:\quad
 L\text{ is related to }U\text{ by finitely many }3\text{-moves and isotopies}.
\]
The statement does not require the unlink to have a prescribed number of components. The supplied research assessment reports that this universal assertion is false: there exists a link that cannot be related by such a sequence to any trivial link. It attributes the negative resolution to Dabkowski--Przytycki (2002), using a Burnside-group obstruction. This is a source-reported resolution; no counterexample or proof is reconstructed here. The generic background note calling the source item open is retained as conflicting status evidence, not used to overwrite the explicit negative assessment.
\subsection{Short English statement}
The conjecture claimed that every link becomes an unlink through finitely many \(3\)-moves. The supplied registry marks it solved negatively and cites a 2002 Burnside-group obstruction. This draft records the original universal claim and its reported disproof, with the local-move convention still requiring source review.
\subsection{Sources}
\begin{itemize}
\item Local UnsolvedMath registry, permanent ID \(10400210\), exact source code \texttt{catalogue entry}, Conjecture 12.4, attributed to Y. Nakanishi [305]. The bibliography entry [305] is absent; the title is truncated.
\item Ohtsuki, \emph{Problems on invariants of knots and 3-manifolds} (2002), Conjecture 12.4, PDF page 160: \url{https://msp.org/gtm/2002/04/gtm-2002-04-024s.pdf}. Bold headings distinguish 231 numbered items from surrounding references.
\item Supplied disproof citation: M. K. Dabkowski and J. H. Przytycki, ``Burnside obstructions to the Montesinos-Nakanishi 3-move conjecture,'' \emph{Geom. Topol.} 6 (2002), 355--360, arXiv:math/0205057. The research summary describes a Burnside-group obstruction.
\item Supplied further references: Dabkowski--Przytycki, ``Unexpected links between components of the same link,'' \emph{Fund. Math.} (2004), for further counterexamples; an earlier Przytycki note (1999) on the \((3,2)\)-move conjecture, without full details. The summary cross-refers to Problem 4.10, source code \texttt{AMR-103-0092}, reporting disproof in February 2002.
\item Extraction pdf-font-structure, accessed 2026-07-29. Registry status solved and classification SOLVED-IN-LITERATURE report a negative resolution, while generic source evidence NEEDS\_REVIEW calls the item open and asks for a current primary-literature audit.
\item Dataset attribution: Ulam AI / UnsolvedMath Contributors, CC BY 4.0. Source rights NEEDS\_REVIEW: publicly accessible author/publisher copy, redistribution terms requiring review.
\end{itemize}
\end{document}
